\usepackage[utf8]{inputenc} % para cargar caracteres no anglosajones 
\usepackage[spanish,english,es-tabla]{babel}
\selectlanguage{spanish}
% Alias para IEEEtran.bst cuando los campos de biblio.bib tienen language = {en}
\makeatletter
\@ifundefined{l@en}{\let\l@en\l@english}{}
\makeatother
\usepackage{datetime} % Opcional, para mejorar el formato de la fecha
\usepackage[lmargin=3cm,rmargin=2.54cm,top=2.54cm,bottom=2.54cm]{geometry}

\usepackage{etoolbox}
%%%%%% PARA UNA FUENTE QUE SE PARECE ARIAL %%%%
%\usepackage[scaled]{helvet}
%\renewcommand\familydefault{\sfdefault}  % Configura Helvetica como la fuente por defecto

\title{Redacción - Tesis - licenciatura}
\author{Juan Alonso Alcalá Lujan}
\date{\today}

% Paquetes simples y estandar 
% % % % % % % % % % % % % % % % % % % % % % % % % % %
\usepackage{bm,amsmath,amssymb,graphicx,amsfonts,amsthm}
\usepackage{fancybox} %este para variar las cajas 
\usepackage{float}
\usepackage{array}
\usepackage{comment} %%% Para comentar secciones largas
%\usepackage[justification=justified]{caption} 
\usepackage[justification=centering, singlelinecheck=false]{caption}
% Luego, reconfigura solo las figuras a justificado
%\captionsetup[figure]{justification=justified}
\usepackage{subcaption} %% para las subfiguras
\usepackage{url}
%\usepackage[numbers]{natbib} % Para usar citas narrativas y paréntetica
\usepackage{lipsum} % Para genera texto random
\usepackage[rgb]{xcolor} % Permite definir colores personalizados en formato RGB
\usepackage{booktabs} %%%%%%% Para las tablas que he exportado en python y Para tablas sin líneas verticales
\usepackage{multicol}
\usepackage{setspace}

% Compatibilidad hacia atrás con tocloft v3.0c en versiones de LaTeX kernel previas a 2024
\providecommand{\UseTaggingSocket}[2][]{}
\providecommand{\SuspendTagging}[1]{}
\providecommand{\ResumeTagging}[1]{}
\usepackage{tocloft}
\usepackage{chngcntr} % Permite cambiar el comportamiento de numeración
\usepackage{appendix}
% Redefiniciones
% % % % % % % % % % % % % % % % % % % % % % % % % % %
\renewcommand{\contentsname}{Tabla de contenidos}
\renewcommand{\bibname}{Bibliografía}
\renewcommand{\indexname}{Índice de Palabras}
\renewcommand{\figurename}{Imagen}
\renewcommand{\tablename}{TABLA}
\renewcommand{\partname}{Parte}
\renewcommand{\chaptername}{Capítulo}
\renewcommand{\appendixname}{Anexos}
\addto\captionsspanish{\renewcommand{\figurename}{FIGURA}}% Cambia "Figura" por "Imagen
\addto\captionsspanish{\renewcommand{\tablename}{TABLA}}

% Definiciones de abreviaturas de revistas astronómicas (usadas en citas BibTeX de NASA/ADS)
\providecommand{\apj}{The Astrophysical Journal}
\providecommand{\apjl}{The Astrophysical Journal Letters}
\providecommand{\apjs}{The Astrophysical Journal Supplement Series}
\providecommand{\aj}{The Astronomical Journal}
\providecommand{\aap}{Astronomy \& Astrophysics}
\providecommand{\mnras}{Monthly Notices of the Royal Astronomical Society}
\providecommand{\pasp}{Publications of the Astronomical Society of the Pacific}
\providecommand{\pasj}{Publications of the Astronomical Society of Japan}
\providecommand{\nat}{Nature}
\providecommand{\sci}{Science}
\providecommand{\araa}{Annual Review of Astronomy and Astrophysics}

% Paquetes rebuscados 
% % % % % % % % % % % % % % % % % % % % % % % % % % %
%%%%%%%%%%%% Para el abstract %%%%%%%%%%%%%%%%%%
\newenvironment{resumen}%
{\null\vfill\begin{center}%
\bfseries \fontsize{14}{16}\selectfont Resumen \end{center}}%
{\vfill\null}

\newenvironment{abstract}%
{\null\vfill\begin{center}%
\bfseries \fontsize{14}{16}\selectfont Abstract \end{center}}%
{\vfill\null}




% COMANDOS 
% % % % % % % % % % % % % % % % % % % % % % % % % % %

\providecommand{\keywords}[1] % Para palabras clave
{
  \small	
  \textbf{\textit{Palabras Clave---}} #1
}



%%%%%%%%% Cosas que no sé para que sirven %%%%%%%%%%%%%
\usepackage{units}
%\usepackage{parskip}


\newcommand{\figref}[1]{Fig. \ref{#1}} % Para modificar la forma como se escribe en la imagenes .
%\usepackage[linktocpage=true,pagebackref = true]{hyperref} 
\usepackage[linktocpage=true]{hyperref} % linktocpage Hace que los números de página en lugar del texto se vinculen en la Tabla de contenido. , %el pagebackref me servira para linkear cosas en citas de bibliografia 


\usepackage{orcidlink}
\definecolor{mycitecolor}{rgb}{0.0, 0.0, 1.0} % Ajusta estos valores para intensificar el color

\hypersetup{ 
    colorlinks=true,
    linkcolor=black, % red
    filecolor=black, % magenta      
    urlcolor=black, % red
    citecolor=black, %mycitecolor        % Color para citas de bibliografía
} %este comando de huperref me sirve para que los limks tengan diferentes colores , depende de lo que necesito

\usepackage{tikz}
\usepackage{empheq}
\usepackage{enumerate}
\usepackage{bm}
\usepackage{physics}
\usepackage{braket} %%%% Paquete que me sirve para dibujar brakets a mi antojo y facilitar cosas
% Para la forma de escribir bra-kets mirar:
% http://mirrors.ibiblio.org/CTAN/macros/latex/contrib/physics/physics.pdf

\usepackage{wasysym}
\usepackage{multirow}
\usepackage{booktabs} % Opcional, mejora la tabla

\renewcommand{\qed}{\hfill$\blacksquare$}

\usepackage{mathrsfs}
\usepackage[makeroom]{cancel}
\usepackage{setspace} % Para establecer espacios específicos

% Keywords comando creado para las palabras clave
\providecommand{\keywords}[1]
{
  \small	
  \textbf{\textit{Palabras Clave---}} #1
}


% Keywords comando creado para las palabras clave
\providecommand{\keywordsenglish}[1]
{
  \small	
  \textbf{\textit{Keywords---}} #1
}




%%%%%%%%%%%%%%%%%%%%%%%%%% SANGRÍA %%%%%%%%%%%%%%%%%%%%%%%%%%%%%%%%%%%%%%%%%%%%%%%%%%%%%%%%%
\usepackage{indentfirst}  % Asegura sangría en el primer párrafo de cada sección
\setlength{\parindent}{1.27cm}  % Define la sangría de todos los párrafos
\linespread{1.25} % Asemeja el interlineado 1.5 de Word

  
%%%%%%%%%%% Para controlar la numeración de pagina %%%%%%%%%%%%%%%%%%%%%%%%%%%%%%%%%%%%%%%%%%%%%%%%%%%%%%%%%%%%%%%%%%%%%%%%%%%%%%%%%%%%%%%%%%%%%%%%%%%%%%%%%%%%%%%%%%%%%
\usepackage{fancyhdr}
\pagestyle{fancy}
\fancyhf{}  % Limpiar encabezados y pies de página
\fancyfoot[C]{} % Borra numeración central
\fancyfoot[R]{\thepage}  % Coloca el número de página en la esquina inferior derecha

% Por defecto, LaTeX usa el estilo plain en las primeras páginas de cada capítulo. Con \fancypagestyle{plain}{...} estamos reconfigurándolo para que las páginas de inicio de capítulo sigan el mismo formato que el resto del documento.
\fancypagestyle{plain}{
    \fancyhf{} 
    \fancyfoot[R]{\thepage} % Número de página en la esquina inferior derecha
}
\renewcommand{\headrulewidth}{0pt}  % Elimina la línea en el encabezado
\renewcommand{\footrulewidth}{0pt}  % Elimina la línea en el pie de página

